%----------------------------------------------------------------------------------------
%    PACKAGES AND THEMES
%----------------------------------------------------------------------------------------

\documentclass[aspectratio=169,xcolor=dvipsnames]{beamer}
\usetheme{SimpleDarkBlue}

\usepackage{hyperref}
\usepackage{graphicx} % Allows including images
\usepackage{booktabs} % Allows the use of \toprule, \midrule and \bottomrule in tables

\usepackage{fontspec}
\usepackage{luatexja}
\usepackage{fontawesome5}

\usepackage{appendixnumberbeamer}

\setmainfont{Alegreya Sans Light}[
  ItalicFont={* Italic},
  BoldFont={Alegreya Sans Medium},
  BoldItalicFont={Alegreya Sans Medium Italic}]
\setsansfont{Alegreya Sans Light}[
  ItalicFont={* Italic},
  BoldFont={Alegreya Sans Medium},
  BoldItalicFont={Alegreya Sans Medium Italic}]

% ------------ %
% beamerbutton %
% ------------ %
\newcommand{\goto}[2]{\hyperlink{#2}{\beamergotobutton{#1}}}
\newcommand{\return}[2]{\hyperlink{#2}{\beamerreturnbutton{#1}}}
\newcommand{\extgoto}[2]{\href{#2}{\beamergotobutton{#1}}}
\newcommand{\orange}[1]{\textcolor{orange}{#1}}
\newcommand{\red}[1]{\textcolor{red}{#1}}
\newcommand{\blue}[1]{\textcolor{blue}{#1}}
\newcommand{\green}[1]{\textcolor{OliveGreen}{#1}}

% % --------------------------- %
% % Section title page with toc %
% % --------------------------- %
% \setbeamertemplate{subsection page}{%
%     \usebeamertemplate*{section page}
% }
% \setbeamertemplate{section in toc}[square]
% \setbeamertemplate{subsection in toc}[square]
% \AtBeginSection[]{
% % \sepframe
% \begin{frame}[noframenumbering]{Outline}
%     % \tableofcontents[currentsection]
%     \tableofcontents[currentsection, currentsubsection]
% \end{frame}
% }
% \AtBeginSubsection[]{
%   \begin{frame}[noframenumbering]{Outline}
%     \tableofcontents[currentsection, currentsubsection]
%   \end{frame}
% }

% ------------------------------------ %
% Section title page with Huge bf text %
% ------------------------------------ %
\AtBeginSection[]{
  \begin{frame}[noframenumbering, plain]
    \Huge{\centerline{\textbf{\insertsection}}}
  \end{frame}
}

%%% automatically add spaces into enumerate and itemize environment
\let\tempone\itemize
\let\temptwo\enditemize
\renewenvironment{itemize}{\tempone\addtolength{\itemsep}{\fill}}{\temptwo}
\let\tempa\enumerate
\let\tempb\endenumerate
\renewenvironment{enumerate}{\tempa\addtolength{\itemsep}{\fill}}{\tempb}

%%=============================================================
%%  FOOTLINE TEMPLATE
%%=============================================================

% % --- Footline depends on onesec option: show only current section or full bar ---
% \if@Optiononesec
% % Minimalist footline: only section name and frame number
% \defbeamertemplate*{footline}{shadow theme}{%
%     \leavevmode\hbox{%
%         \hypersetup{linkcolor=white,urlcolor=white,citecolor=white}
%         \begin{beamercolorbox}[wd=1.01\paperwidth, ht=2.5ex, dp=1.125ex, leftskip=.3cm plus1fil, rightskip=0.3cm]{author in head/foot}%
%             \insertshorttitle \insertsection \hfill \insertframenumber \ / \ \inserttotalframenumber%
%         \end{beamercolorbox}%
%     }%
% }
% \else
% % Standard footline: section bar and frame number
\defbeamertemplate*{footline}{shadow theme}{%
    \leavevmode\hbox{%
        \hypersetup{linkcolor=white,urlcolor=white,citecolor=white}
        \begin{beamercolorbox}[wd=1.02\paperwidth, ht=2.5ex, dp=1.125ex, leftskip=.3cm plus1fil, rightskip=0.3cm]{author in head/foot}%
            \insertsectionnavigationhorizontal{.80\textwidth}{}{} \hspace{.3cm} \hfill \#\ \insertframenumber \ / \ \inserttotalframenumber%
        \end{beamercolorbox}%
    }%
}
% \fi

% --- Activate the custom footline ---
\setbeamertemplate{footline}[shadow theme]


%----------------------------------------------------------------------------------------
%    TITLE PAGE
%----------------------------------------------------------------------------------------

\title[Asset Pricing (Production)]{Asset Pricing in Production Economy}
\author[Hui-Jun Chen]{Hui-Jun Chen}
\institute[NTHU]{National Tsing Hua University}
\vspace{2pt}
\date{\today}

\begin{document}

%----------------------------------------------------------------------------------------
%    PRESENTATION SLIDES
%----------------------------------------------------------------------------------------

\begin{frame}{Overview}
\label{slide:Overview}
    \begin{center}
        How does share price comove with GDP?
    \end{center}
    \begin{itemize}
        \item We extend Lucas (1978) to production economy $ \Rightarrow  $ \textbf{firms}
        \item firms are \blue{active} player in macro: \textbf{investment v.s. GDP volatility}
        \begin{itemize}
            \item corporate finance: firm debt? capital investment?
            \item human resource: hiring / lay off employee?
            \item international economics: multi-nation enterprise? FDI?
        \end{itemize}
        \item To be able to reach some conclusion, we need simplification:
        \begin{itemize}
            \item similar setting as Lucas (1978), representative HH \& firm
            \item firm pays dividend $ \Leftarrow  $ firm are \textbf{DRS}
            \item labor-only technology $ \Rightarrow  $ no other intertemporal asset other than share.
        \end{itemize}
    \end{itemize}
\end{frame}

\section[Firm]{Firm Problem}
\label{sec:Firm_Problem}

\begin{frame}{Dividend and Wage}
\label{slide:Dividend_and_Wage}
\begin{itemize}
    \item Production function: $ \displaystyle y = z n^{\alpha} $, where $ z $ is TFP shock, and $ \alpha \in (0, 1) $.
    \item Firm's profit maximization problem: $ \displaystyle \max_{n} z n^{\alpha} - wn $
    \begin{itemize}
        \item FOC: $ \displaystyle w = \alpha z n ^{\alpha - 1} $
    \end{itemize}
    \item Wage bill: $ \displaystyle w n = \alpha z n ^{\alpha} = \alpha y $
    \item Assume firm all profits as dividend, $ \displaystyle d = y - wn = (1-\alpha)y $
\end{itemize}

\end{frame}

\section[Household]{Household's Problem}
\label{sec:Household_s_Problem}

\begin{frame}{Household Problem}
\label{slide:Household_Problem}

Assume HH value leisure, and thus
%
\begin{align}
    V(s, z)
        & = \max_{c\ge 0, s'\ge 0, n \ge 0} \log c + \psi (1-n) + \beta \mathbb{E}_{z'|z} [V(s', z')]
    \\
    \text{s.t. } \quad
        & c + p s' \le (d + p) s + wn
\end{align}
%
We know in equilibrium / steady state, three markets need to clear:
\begin{enumerate}
    \item find $ w $ such that labor demand $ = $ labor supply
    \item find $ p $ such that $ s = 1 $
    \item by Walras' law, goods market clear, implying $ c = y $.
\end{enumerate}

\end{frame}

\begin{frame}{Solve Household Problem}
\label{slide:Solve_Household_Problem}
    Using the same solution technique,
    %
    \begin{align}
        V(s, z)
            & = \max_{s', c, n} \log c + \psi (1-n) + \beta \mathbb{E}_{z'|z}[\log c' + \psi (1-n')]
        \\
            & \quad + \beta^{2} \mathbb{E}_{z'|z} [V(s'', z'')]
        \\
        \text{subject to } \quad
            & c + ps' \le (d + p)s + wn
        \\
            & c' + p's'' \le (d' + p') s' + w' n'
    \end{align}
    %
    Replace $ c $ and $ c' $ and get
    %
    \begin{align}
        V(s, z)
            & = \max_{s', n} \log((d+p)s + wn - p s') + \psi (1-n)
        \\
            & \quad + \beta \mathbb{E}_{z'|z} [\log((d'+p')s' + w'n' - p's'') + \psi (1 - n')]
        \\
            & \quad + \beta^{2} \mathbb{E}_{z'|z} [V(s'', z'')]
    \end{align}
    %
\end{frame}

\begin{frame}{First Order Condition}
\label{slide:First_Order_Condition}
    %
    \begin{align}
        V(s, z)
            & = \max_{s', n} \log((d+p)s + wn - p s') + \psi (1-n)
        \\
            & \quad + \beta \mathbb{E}_{z'|z} [\log((d'+p')s' + w'n' - p's'') + \psi (1 - n')]
        \\
            & \quad + \beta^{2} \mathbb{E}_{z'|z} [V(s'', z'')]
    \end{align}
    %
    FOC:
    %
    \begin{align}
        [n]: \quad
            & \frac{w}{c} = \psi
        \\
        [s']: \quad
            & \frac{1}{c} \cdot p = \beta \mathbb{E}_{z'|z} \left[
                \frac{1}{c'} \cdot (d' + p')
            \right]
    \end{align}
    %
\end{frame}

\section[Eqm]{Equilibrium Outcome}
\label{sec:Equilibrium_Outcome}

\begin{frame}{Optimality Conditions}
\label{slide:Optimality_Conditions}
    \begin{align}
        \label{eq:NFOC}
        [n]: \quad
            & \frac{w}{c} = \psi \Rightarrow w = \psi c
        \\
        \label{eq:SprimeFOC}
        [s']: \quad
            & \frac{1}{c} \cdot p = \beta \mathbb{E}_{z'|z} \left[
                \frac{1}{c'} \cdot (d' + p')
            \right]
        \\
        \label{eq:NDFOC}
        [\text{Firm}]: \quad
            & w = \alpha z n^{\alpha-1}
    \end{align}
    $ w = w $, \eqref{eq:NFOC} equals to \eqref{eq:NDFOC}, and $ c = y $ yields
    %
    \begin{align}
        \psi y = \alpha \frac{y}{n}
            & \Rightarrow n = \frac{\alpha}{\psi} \Rightarrow y = z n^{\alpha} = z \left(
                \frac{\alpha}{\psi}
            \right)^{\alpha}
        \\
            & \Rightarrow w = \alpha z n ^{\alpha-1} = \alpha z \left(
                \frac{\alpha}{\psi}
            \right)^{\alpha-1}
        \\
            & \Rightarrow d = (1-\alpha) y = (1-\alpha) z \left(
                \frac{\alpha}{\psi}
            \right)^{\alpha}
    \end{align}
    %

\end{frame}

\begin{frame}{Share Euler Equation}
\label{slide:Share_Euler_Equation}
    Focus on \eqref{eq:SprimeFOC}, we can use $ c' = y' $ as well as $ d' = (1-\alpha) y' $ to simplify:
    %
    \begin{align}
        \frac{p}{y}
            & = \beta \mathbb{E}_{z'|z} \left[
                \frac{p'}{y'} + \frac{(1-\alpha)y'}{y'}
            \right]
        \\
            & = \beta (1-\alpha) + \beta \mathbb{E}_{z'|z} \left[
                \frac{p'}{y'}
            \right]
    \end{align}
    %
    Somehow you got a prophecy from the spirit and his/her voice tells you to guess $ \displaystyle \frac{p}{y} \equiv \Lambda $, a constant over time regardless of TFP shock. Is that true?
    %
    \begin{align}
        \Lambda
            & = \beta (1-\alpha) + \beta \mathbb{E}_{z'|z} \left[
                \Lambda
            \right] = \beta(1-\alpha) + \beta \Lambda
        \\
        \Lambda
            & = \frac{\beta(1-\alpha)}{1-\beta}
    \end{align}
    %
    \begin{center}
        It true \faAngellist \faAngellist \faAngellist
    \end{center}
\end{frame}

\begin{frame}{Intepretation}
\label{slide:Intepretation}
    Stock price to GDP ratio, $ \frac{p}{y} $, is constant over time, which implies
    \begin{enumerate}
        \item stock price is procyclical: \faArrowCircleUp \hspace{1px} and \faArrowCircleDown \hspace{1px} with TFP $ z $,
        \item the percentage std of stock price matches percentage std of dividend,
        \item stock is risky: $ \displaystyle p = \frac{\beta(1-\alpha)}{1-\beta} y $ $ \Rightarrow  $ requires $ (+) $ risk premium
        \begin{itemize}
            \item $ \displaystyle e(z, z') = \frac{d'+p'}{p} = \frac{(1-\alpha)y' + \Lambda y'}{\Lambda y} = \frac{\frac{1-\alpha}{1-\beta}y'}{ \frac{\beta(1-\alpha)}{1-\beta} y} = \frac{1}{\beta} \frac{y'}{y}$
            \item SDF $ \displaystyle = \frac{\beta u'(c')}{u'(c)} = \beta \frac{y}{y'} $
            \item Risk premium $ \displaystyle = \frac{\mathbb{E}_{t}[e(z, z') - R_{t}]}{R_{t}} = - cov_{t} \left[
                SDF, e(z, z')
            \right] > 0$
        \end{itemize}
    \end{enumerate}
    The very times firm shares pay high is when your consumption is low!
\end{frame}

%\section{Appendix}
%\label{sec:Appendix}

%\appendix
%% -------------------------------------------
%\setbeamertemplate{headline}
%{
%\setbeamercolor{section in head/foot}{fg=black, bg=white}
%\vskip1em \tiny \insertsectionnavigationhorizontal{1\paperwidth}{\hspace{0.50\paperwidth}}{}
%}
%%------------------------------------------
%% \begin{frame}\frametitle{}
%% \begin{columns}
%% \label{Appendix}
%% \column{1\linewidth}
%% \centering
%% {\Large \alert{Appendix}}
%% \end{columns}
%% \end{frame}
%%------------------------------------------
%\begin{frame}[allowframebreaks]{References}
%\footnotesize
%\bibliographystyle{$BIB_STYLE}
%\bibliography{$BIBFILE}
%\end{frame}

\end{document}
